\documentclass[12pt,a4paper]{article}
\usepackage[margin=2.5cm]{geometry}
\usepackage{amsmath,amssymb,amsfonts}
\usepackage{booktabs}
\usepackage{xcolor}
\usepackage{graphicx}
\usepackage{hyperref}
\usepackage{caption}
\usepackage{array}
\usepackage{enumitem}
\usepackage{fancyhdr}

\hypersetup{colorlinks=true, linkcolor=blue, citecolor=blue, urlcolor=blue}
\pagestyle{fancy}
\fancyhf{}
\rhead{QML Audit Report --- Paper 02}
\lhead{Autonomous ML Research Review}
\cfoot{\thepage}

\title{\textbf{Audit \& Reproduction Report}\\[0.4em]
\large\textit{Paper 02: Quantum-Inspired Stacked Integrated Concept Graph Model (QISICGM) for Diabetes Risk Prediction}\\[0.3em]
Young, arXiv:2509.12259 (2025)}
\author{Autonomous Senior ML Research Engineer}
\date{\today}

\begin{document}
\maketitle
\tableofcontents
\newpage

%---------------------------------------------------------------
\section{Paper Overview}
%---------------------------------------------------------------
\subsection{Bibliographic Details}
\begin{itemize}[leftmargin=*]
  \item \textbf{Author:} Kevin Young
  \item \textbf{Venue:} arXiv preprint, arXiv:2509.12259 (2025)
  \item \textbf{Repository:} \url{https://github.com/keninayoung/QISICGM}
  \item \textbf{Topic:} ``Quantum-Inspired'' stacked ensemble with concept graph construction for diabetes risk prediction on Pima dataset.
\end{itemize}

\subsection{Goal}
Introduce QISICGM, a hybrid classifier combining:
(a) \emph{quantum-inspired} phase feature maps (trigonometric feature encoding),
(b) a kNN-based concept graph for neighbourhood relationship modelling,
(c) a stacked ensemble of Random Forest, Extra Trees, and a Feed-Forward Neural Network (FFNN) base learners, and
(d) a Logistic Regression meta-learner trained on out-of-fold (OOF) predictions.

\subsection{Dataset}
\begin{itemize}[leftmargin=*]
  \item \textbf{Pima Indians Diabetes:} 768 rows, 8 continuous features.
  \item \textbf{Synthetic minority augmentation:} 2,000 additional synthetic minority-class samples generated via Gaussian Mixture Models (GMM), fitted to the \emph{full} Pima training data --- constituting a data-leakage risk.
\end{itemize}

\subsection{Quantum-Inspired Feature Map}
For each feature $x_k$, the QISICGM phase map computes:
\begin{equation}
  \phi(x_k) = \bigl[\cos(\alpha x_k),\; \sin(\alpha x_k)\bigr]
\end{equation}
where $\alpha$ is a scaling factor. This is a \emph{purely classical} trigonometric transformation, not a quantum computation, despite the paper's quantum framing.

\subsection{Architecture}
\begin{enumerate}[leftmargin=*]
  \item \textbf{Phase Feature Map:} $8 \to 16$ dimensional cos/sin expansion.
  \item \textbf{Concept Graph:} kNN ($k=5$) adjacency, with node features propagated via mean aggregation.
  \item \textbf{Base Learners:} RF (100 trees), Extra Trees (100 trees), FFNN (3 hidden layers).
  \item \textbf{Meta-Learner:} LogisticRegression on OOF predictions, 5-fold OOF.
  \item \textbf{Calibration:} Isotonic regression calibration applied to base learner probabilities.
\end{enumerate}

\subsection{Claimed Performance}
\begin{itemize}[leftmargin=*]
  \item Accuracy: \textbf{99.8\%}
  \item ROC-AUC: \textbf{0.999}
  \item F1-Score: \textbf{99.7\%}
  \item Brier Score: not reported
\end{itemize}

%---------------------------------------------------------------
\section{Red Flags \& Methodological Audit}
%---------------------------------------------------------------
\begin{enumerate}[leftmargin=*]
  \item \textbf{GMM Synthetic Leakage (Critical):} The 2,000 synthetic minority samples are generated by fitting GMM on the full Pima dataset (including rows later designated as validation/test). When these synthetic rows appear in the training fold but their generating GMM was fitted on test rows, performance is artificially inflated.
  \item \textbf{``Quantum-Inspired'' is purely classical:} The phase map $[\cos(\alpha x), \sin(\alpha x)]$ is a standard Fourier feature expansion (Rahimi \& Recht, 2007). No quantum circuits, qubits, or quantum operations are used. The ``quantum'' branding is misleading.
  \item \textbf{Synthetic rows inside CV folds:} Even when using OOF evaluation, synthetic rows that are statistically derived from test-fold rows inflate OOF metrics. Clean evaluation requires: (a) fit GMM on training split only, (b) augment training split only, (c) evaluate meta-learner on real rows only.
  \item \textbf{Accuracy of 99.8\% on Pima:} Well-known to be impossible with honest evaluation on this dataset (Bayes ceiling $\approx 78$--80\%). All published leakage-free benchmarks report 74--80\%.
\end{enumerate}

%---------------------------------------------------------------
\section{Faithful Reimplementation Details}
%---------------------------------------------------------------
Implemented in \texttt{notebooks/02\_young\_qisicgm.ipynb}:
\begin{itemize}[leftmargin=*]
  \item Phase feature map ($\cos/\sin$) applied to 8 features $\to$ 16 dims.
  \item kNN concept graph ($k=5$) with mean propagation.
  \item RF + ET + FFNN base learners with 5-fold OOF meta-features.
  \item LogisticRegression meta-learner.
  \item \textbf{Scenario A (Leaky):} GMM fitted on full 768-row dataset before fold generation. Reflects paper's procedure.
  \item \textbf{Scenario B (Clean):} GMM fitted \emph{inside each training fold only}. Meta-learner evaluated on real rows only (no synthetic rows in test/validation).
\end{itemize}

%---------------------------------------------------------------
\section{Results: Paper-Claimed vs.\ Reproduced vs.\ Leakage-Free}
%---------------------------------------------------------------
\begin{table}[h!]
  \centering
  \caption{QISICGM: Claimed vs.\ Reproduced vs.\ Leakage-Free (Pima dataset)}
  \begin{tabular}{lccc}
    \toprule
    \textbf{Metric} & \textbf{Paper Claimed} & \textbf{Reproduced (Leaky)} & \textbf{Clean Real Rows}\\
    \midrule
    Accuracy   & 99.8\%  & 90.75\% & 75.91\%\\
    ROC-AUC    & 0.999   & 0.9631  & 0.8284\\
    F1-Score   & 99.7\%  & 93.42\% & 61.86\%\\
    Brier Score& N/R     & 0.0685  & 0.1600\\
    \bottomrule
  \end{tabular}
\end{table}

\subsection{Quantified Inflation}
\begin{itemize}[leftmargin=*]
  \item \textbf{Accuracy inflation:} 99.8\% $\to$ 75.9\% = \textcolor{red}{\textbf{$-23.9$ pp}} when leakage removed.
  \item \textbf{GMM leakage contribution:} Reproduced (leaky) 90.75\% $\to$ Clean 75.91\% = \textcolor{red}{\textbf{$+14.84$ pp}} inflation from synthetic augmentation leakage.
  \item \textbf{F1 inflation:} 99.7\% $\to$ 61.86\% = \textcolor{red}{\textbf{$+31.56$ pp}} inflation; class imbalance on Pima makes F1 particularly sensitive to leaky oversampling.
\end{itemize}

%---------------------------------------------------------------
\section{Discussion}
%---------------------------------------------------------------
The QISICGM's real contribution is a well-constructed stacked ensemble with a kNN concept graph. When evaluated correctly (no synthetic leakage, no quantum branding), it achieves a ROC-AUC of 0.828 and accuracy of 75.9\% on Pima --- statistically competitive with classical RBF-SVM but not superior.

The ``quantum-inspired'' framing via $\cos/\sin$ features adds no demonstrated benefit beyond standard Fourier feature expansion and could mislead readers into conflating trigonometric preprocessing with quantum computation.

%---------------------------------------------------------------
\section{Conclusion}
%---------------------------------------------------------------
QISICGM's claimed 99.8\% accuracy is entirely attributable to GMM-synthetic data leakage. Clean evaluation yields a competitive but ordinary 75.9\% accuracy, indistinguishable from a well-tuned Random Forest on Pima Indians Diabetes. The quantum-inspired branding is scientifically unsupported.

\appendix
\section{Artifacts}
\begin{itemize}
  \item Notebook: \texttt{notebooks/02\_young\_qisicgm.ipynb}
  \item Metrics CSV: \texttt{results/paper02\_young\_metrics.csv}
  \item Figure: \texttt{results/fig\_paper02\_calibration.png}
\end{itemize}

\end{document}
